\appendix
\section{Survey methods and corpus construction}\label{app:methods}

The headline claims of this survey, in particular the lifecycle-stage and
state-axis counts and the rollback and authority figures the argument
rests on, are only as credible as the procedure that produced the corpus. This
appendix documents that procedure so the numbers are auditable rather than
asserted.

\subsection{Sources and time window}

The corpus draws on four sources: arXiv (cs.AI, cs.CL, cs.LG, cs.CR) as the
primary source for 2023 to 2026 preprints; Semantic Scholar for backward and
forward citation chasing and for venue and citation metadata; OpenReview (ICLR,
NeurIPS, COLM) for venue status and camera-ready versions; and the ACL Anthology
for published versions at ACL, EMNLP, and NAACL. The time window is 2023-01
through 2026-06, with foundational anchors admitted regardless of date when they
ground a taxonomy or a term the survey uses (for example, the episodic-semantic
distinction, catastrophic interference, and the Soar architecture).

\subsection{Search procedure}

Collection ran in three passes. We \emph{seeded} from recent agent-memory surveys
and the closest-work set and extracted their unit of analysis and shared
vocabulary. Next we \emph{chased citations} backward (references) and forward
(citing papers via Semantic Scholar) two hops from each seed, keeping any work
that touches a persistent-state operation, state type, benchmark, or failure
mode. Finally we ran \emph{targeted term-set sweeps} for 2024 to 2026 work the
citation graph missed, with query sets grouped by mechanism (for example
``agent memory'', ``memory OS'', ``procedural memory agent'', ``knowledge graph
memory''), benchmark (``long-term conversational memory'', ``multi-session
benchmark'', ``lifelong agent benchmark''), failure mode (``memory poisoning LLM
agent'', ``indirect prompt injection'', ``machine unlearning LLM'', ``deletion
residue''), continual adaptation and personalization (``continual learning
agent'', ``test-time training'', ``knowledge editing forgetting''), and
foundations (``cognitive architecture memory'', ``LLM agent survey''). The
targeted sweeps were run as iterative non-duplicate rounds: each round queried
the sources, removed works already in the corpus, and added only genuinely new
hits, until additional rounds returned predominantly duplicates. Rounds aimed
deliberately at the governance frontier (rollback, authority, deletion
propagation, durable execution, formal verification, and the adjacent database,
distributed-systems, and HCI literatures) so that the governance counts cannot be
attributed to having searched only the memory mainstream.

\subsection{Inclusion and exclusion criteria}

A work is \emph{included} if it satisfies at least one of: (I1) it builds or
studies a persistent-state mechanism; (I2) it contributes a benchmark or
evaluation method relevant to memory, long-horizon, or stateful behavior; (I3) it
documents a persistence-induced failure mode; (I4) it is a foundational anchor a
reader needs to place the state taxonomy, the lifecycle, or the
plasticity-stability framing; or (I5) it defines a boundary case (long-context,
RAG, personalization) used to delimit always-on agents. A work is \emph{excluded}
if it is (E1) pure model-architecture or pretraining work with no memory, state,
or agent angle; (E2) an application that uses an agent but contributes nothing to
a persistent-state operation, state type, benchmark, or failure mode (we add no
padding for breadth alone); (E3) a duplicate or superseded version when a later
one is cited; or (E4) available only in a venue without an accessible English
record. The resulting corpus is $N{=}435$ works.

\subsection{Coding scheme}

Every included work is coded along four dimensions: a \emph{category} (mechanism,
benchmark, failure mode, foundation, boundary, or survey); the \emph{lifecycle
stages} it exercises, as a multi-label set over observe, write, validate,
organize, retrieve, act, update, forget, audit, and rollback; the \emph{state
axes} it addresses, as a multi-label set over authority, scope, mutability,
provenance, recoverability, and actionability; and an \emph{application sub-area}.
Multi-label coding is why per-stage and per-axis incidences sum to more than the
corpus size. All counts in the survey are recomputed from the coded corpus by a
script rather than hand-entered.

\subsection{Coding reliability}

Because the lifecycle and axis labels are interpretive, we measured inter-coder
reliability with a blind second coding of a $236$-work sample (over half the
corpus). On this sample, pooled per-cell agreement is $0.82$ on lifecycle stages
(Cohen's $\kappa = 0.58$) and $0.74$ on state axes ($\kappa = 0.44$), which is
moderate agreement, as expected for fine-grained multi-label coding of short
descriptions. Two points matter for how the headline claims should be read.
First, the reliability evidence supports the survey's \emph{directional} claims
(forward arc dominant, return arc sparse) far better than it supports any exact
cell count, and that is the only weight we place on it. Even the rollback label,
the rarest, carries real labeling noise: the two coders agreed on it for only
$11$ of the $19$ works either marked, so the precise value $27$ should be read as
an estimate with a margin of several works, not a measurement. Second, the
directional comparison does not depend on that precision, because it is a large
qualitative asymmetry rather than a narrow margin: retrieve appears in
$269$ works and rollback in roughly $27$, an order-of-magnitude gap that no
plausible reclassification of borderline cells could close. A reader who
distrusts the exact $27$ can inflate it severalfold and still obtain the same
directional conclusion: the return arc is much sparser than the forward arc. We
therefore treat the counts as well-supported for the comparative claims the
survey makes (forward arc dominant, return arc sparse) while cautioning that any
single cell
should be read as an estimate, not a precise measurement.

\subsection{What the rollback count does and does not aggregate}

The figure that twenty-seven of $435$ works expose any rollback mechanism
aggregates three mechanically distinct operations: internal-state rollback
(reverting a stored fact or a consolidation decision), workflow rollback
(compensating a partially executed multi-step task), and external-effect rollback
(undoing an already-committed side effect such as a payment or a deletion). The
aggregate is therefore generous for the easiest subclass and an overcount for the
hardest. External-effect rollback, the subclass that matters most because
irreversible effects cannot be repaired by any internal mechanism, is rarer than
the aggregate suggests. A finer corpus pass should separate the three; we report
the conservative aggregate and flag the subclass structure so the headline number
is not over-read.

\subsection{Threats to validity}

Three threats deserve naming. First, the corpus is a \emph{representative coding,
not an exhaustive census}; we sampled toward the governance frontier and the gap
persisted, but we cannot exclude that some governance work uses vocabulary none of
our query sets reached, which is why the universal claims in the body are scoped
to the corpus or to works we examined rather than to the literature as a whole.
Second, the corpus is weighted toward recent preprints, and the governance turn we
document is itself recent (Figure~\ref{fig:timeline}); restricting to work
through 2024 makes the governance fraction smaller rather than larger, so the gap
is not an artifact of including immature 2026 work, but the durability of the
2025 to 2026 governance preprints remains unproven. Third, single-coder labeling
with a second-coder check is weaker than full double-coding; the $\kappa$ values
above bound how much this matters, but they do not make the exact cell counts
immune to recoding.
